\section{Compiling games with the data hiding test}\label{sec:compile-sec}

Now we can show how to compile games from the second layer to the first layer.
Our construction is given in the following definition.

\begin{definition}
Let $\game_k$ be a $(k, n, q)$-register game.
Then its compiled version is the game $\calC_{k \rightarrow \mathrm{semi}}(\game_k)$ defined in \Cref{fig:two-to-one}.
\end{definition}
{
\floatstyle{boxed} 
\restylefloat{figure}
\begin{figure}
With probability~$\tfrac{1}{2}$ each, perform one of the following three tests.
\begin{enumerate}
	\item \textbf{Data hiding:} Draw $(\bx, \bx', \bC) \sim \game_k$, where $\bx = (\bx_1, \bx_2)$ and $\bx_1 = (\bW_1, \ldots, \bW_k)$.
			 If $\bW_k = \hideq$, play $\game_{\mathrm{hide}}$ with question~$\bx$.\label{item:data-hiding}
	\item \textbf{Play game:} Perform $\game_{k}$.\label{item:play-ball!}
	\end{enumerate}
	\caption{The game $\calC_{k \rightarrow \mathrm{semi}}(\game_k)$.\label{fig:two-to-one}}
\end{figure}
}

\begin{theorem}\label{theorem:two-to-one-compiler}
Suppose $\game_k$ is a $(k, n, q)$-register game,
and consider the $(k, n, q)$-semiregister game $\game_{\mathrm{semi}} = \calC_{k \rightarrow \mathrm{semi}}(\game_k)$.
\begin{itemize}
\item[$\circ$] \textbf{Completeness:} 
	Suppose there is a value-$1$ $(k,n,q)$-register strategy for $\game_{k}$ which is also a real commuting EPR strategy.
	Then there is a value-$1$ $(k,n,q)$-semiregister strategy for $\game_{\mathrm{semi}}$ which is also a real commuting EPR strategy.
\item[$\circ$] \textbf{Soundness:} If $\valsemi{k,n,q}{\game_{\mathrm{semi}}} \geq 1-\eps$ then $\valreg{k,n,q}{\game_{k}} \geq 1 -\delta(\eps)$, where $\delta(\eps) = \mathrm{poly}(\eps)$.
\end{itemize}
Furthermore,
\begin{equation*}
\qlength{\game_{\mathrm{semi}}} = O(\qlength{\game_k}), \quad
\alength{\game_{\mathrm{semi}}} = O(\alength{\game_k}),
\end{equation*}
\begin{equation*}
\qtime{\game_{\mathrm{semi}}} = O(\qtime{\game_k}), \quad
\atime{\game_{\mathrm{semi}}} = O(\atime{\game_k}).
\end{equation*}
\end{theorem}

\newcommand{\phoenicianindahizzous}{\textphnc{h}}

\begin{proof}[Proof of \Cref{theorem:two-to-one-compiler}]
The communication and time complexities are the result of combining
the communication and time complexities from $\game_{\mathrm{k}}$ with the values for the data hiding game from \Cref{def:data-hide-def}.

\paragraph{Completeness.}
Let $(\psi, M)$ be a value-$1$ $(k,n,q)$-register strategy for $\game_{k}$ which is also a commuting EPR strategy.
Then for every $x = (x_1, x_2)$ where $x_1 = (W_1, \ldots, W_k)$ with $W_k = \hideq$,
by \Cref{thm:data-hiding-game}
we can extend this strategy to one that passes the data hiding game with question~$x$ with probability~$1$.
Thus, this strategy has value~$1$ overall.
In addition, \Cref{thm:data-hiding-game} implies this strategy is a real commuting EPR strategy as well.

\paragraph{Soundness.}
Suppose $\calS = (\psi, M)$ is a $(k, n, q)$-semiregister strategy for $\game_{\mathrm{semi}}$ with value $1-\eps$.
By \Cref{lem:proj-suffices}, we can assume without loss of generality that~$M$ is projective.
Our goal will be to decode $\calS$ into a $(k, n, q)$-register strategy $\calS_{k}$ with nearly the same value.

\paragraph{Using the data hiding test.}
For a fixed question~$x$, write $\nu_x$ for the probability that~$\calS$ passes the test in \Cref{item:data-hiding}.
Then on average, the probability that~$\calS$ passes this test is $\E_{\bx} \nu_{\bx}$,
which is at least $1-2\eps$ because the overall test passes with probability at least $1-\eps$.
This implies that $\nu_{\bx} \geq 1-\eps^{1/2}$ with probability at least $1-2\eps^{1/2}$.
Given a matrix~$M$ and a $W \in \{X, Z, \hideq, \noop\}$, let us write $\hide{W}{M}:= \hide{k}{M}$ if $W = \hideq$
and $\hide{W}{M}:=M$ otherwise.
For a question~$x$, if $W_k \neq \hideq$, then $\hide{W_k}{M^x_a} = M^x_a$ trivially.
On the other hand, suppose $W_k = \hideq$.
Then either $\nu_x \geq 1-\eps^{1/2}$,
in which case $M^x_a \otimes I_{\reg{Bob}} \approx_{\delta(\eps)} \hide{W_k}{M^x_a}\otimes I_{\reg{Bob}}$ by \Cref{thm:data-hiding-game},
or $\nu_x < 1-\eps^{1/2}$,
in which case we have the trivial bound $M^x_a \otimes I_{\reg{Bob}} \approx_{1} \hide{W_k}{M^x_a}\otimes I_{\reg{Bob}}$ from~\Cref{fact:trivial-upper-bound-approx-delta}.
Since this latter case happens with probability at most $2\eps^{1/2}$,
averaging over all~$x$ gives us
\begin{equation}\label{eq:dr-jekyll-mr-hide}
M^x_a \otimes I_{\reg{Bob}} \approx_{\delta(\eps)} \hide{W_k}{M^x_a}\otimes I_{\reg{Bob}},
\end{equation}
on the distribution $\calD$.

\paragraph{Extracting a strategy.}
Define the strategy $\calS_k = (\psi, \Lambda)$, in which
$\Lambda_a^x := \hide{W_k}{M^x_a}$.
First, we show that $\calS_k$ is a $(k, n, q)$-register strategy.
To do so, fix $x= (x_1, x_2)$ with $x_1 = (W_1, \ldots, W_k)$ and $a = (a_1, a_2)$ with $a_1 = (u_1, \ldots, u_k)$.
Then
\begin{equation*}
\Lambda_{a_1}^x
= \mathrm{hide}_{W_k}\big(\tau^{W_1}_{u_1} \otimes \cdots \otimes \tau^{W_k}_{u_k} \otimes I_{\reg{aux}}\big)\\
	= \tau^{W_1}_{u_1} \otimes \cdots \otimes \tau^{W_k}_{u_k} \otimes I_{\reg{aux}}.
\end{equation*}
The first equality is by definition of~$\Lambda$ and the fact that $\calS$ is a $(k, n, q)$-quasiregister strategy.
The second equality is trivial when $W_k \neq \hideq$ and follows from the fact that $\tau^{W_k}_{\varnothing} = I$ when $W_k = \hideq$.
Next, define $S = \{i\neq k  \mid W_i = \hideq\}$.
If $W_k \neq \hideq$ then $\Lambda_a^x = M^x_a = M_{\overline{S}} \otimes I_S$ for some matrix~$M$.
Otherwise, if $W_k = \hideq$, set $\phoenicianindahizzous = \tr(I_k)$. Then
\begin{equation*}
\Lambda_a^x = \hide{k}{M^x_a} = \hide{k}{M_{\overline{S}} \otimes I_S}
= \frac{1}{\phoenicianindahizzous} \cdot I_k \otimes \tr_k(M_{\overline{S}} \otimes I_S)
= \frac{1}{\phoenicianindahizzous} \cdot I_{S \cup k} \otimes \tr_k(M_{\overline{S}}).
\end{equation*}
The matrix $\tr_k(M_{\overline{S}}) \cdot \phoenicianindahizzous^{-1}$
only acts on the registers in $\overline{S \cup k}$ and the auxiliary register,
and as a result, this strategy satisfies data hiding.
Thus, $\calS_k$ is a $(k, n, q)$-register strategy.

It remains to show that $\calS_k$ has good value.
This follows by combining~\Cref{eq:dr-jekyll-mr-hide} with~\Cref{fact:approx-delta-game-value}:
$\valstrat{\game_{k}}{\calS_{k}} \geq \valstrat{\game_{\mathrm{semi}}}{\calS} - \delta(\eps)$,
and so $\valreg{k,n,q}{\game_{k}} \geq 1 - \delta(\eps)$.
\end{proof}

%%% Local Variables:
%%% mode: latex
%%% TeX-master: "main"
%%% End:
%%%---------- close: data-hiding-compiler

%%%%%%%%%%% jump to partial-data-hiding
%%%---------- open: partial-data-hiding
%!TEX root = main.tex


